\section{端到端配方：从 OLMo 到 Tulu 的组合逻辑}

\subsection{组合原则：先稳底座，再做对齐，最后做预算优化}
这场讲座给出的真正实践路线不是 ``先选算法''，而是 ``先选组合顺序''。一个可执行的顺序是：
\begin{enumerate}
    \item 预训练阶段先建立稳定底座；
    \item SFT 把任务协议和输出行为校正到可控范围；
    \item DPO/PPO 做偏好细调；
    \item RLVR 在可验证子域冲刺推理；
    \item 测试时扩展服务于高价值长尾请求。
\end{enumerate}

\begin{knowledgebox}{为什么这一路线对开放社区友好}
每一层都可以独立实验、独立评估、独立开源，降低了研究协作门槛，也让复现实验失败时更容易定位问题所在。
\end{knowledgebox}

\subsection{评估体系：不要只看一个 benchmark}
\begin{importantbox}{讲座中的评估观}
Reasoning 提升必须跨多个维度共同验证：
\begin{itemize}
    \item 数学/代码等可验证任务；
    \item 知识问答与事实性；
    \item 长上下文稳定性；
    \item 多语言与安全性约束；
    \item 成本与延迟指标。
\end{itemize}
\end{importantbox}

\begin{warningbox}{指标成功不等于系统成功}
如果评估不覆盖部署约束（延迟、预算、鲁棒性），模型即使 benchmark 提升，也可能在真实产品环境不可用。
\end{warningbox}

\subsection{本章小结}
端到端配方的核心是顺序和组合，而不是单点最优。成功系统通常是多阶段 ``次优但互补'' 的结果。


